\section{External Validation: Expert Review Protocol}
\label{sec:expertreview}

The audit in Section~\ref{sec:judgeaudit} compares rubrics that share a
model substrate. Breaking out of that loop requires a judgment the loop
cannot produce, which is the judgment of a human regulatory professional.
I describe the protocol here, in advance of its results, so that the
design of the external validation is on the record and does not depend on
how the results turn out. Labels are being collected at the time of
writing. What follows is the sampling frame, the measurement plan, and
the outcomes that would count against the claims made in this paper.

\subsection{Sampling frame}
I drew a stratified sample of 50 items from the frozen bank, ten from
each of the five hard domains. Items were restricted to those passing a
tightened quality filter (point-biserial $\geq 0.3$, infit $\leq 1.3$,
outfit $\leq 1.5$) and to the interquintile difficulty band
$b \in [0.83, 2.43]$, which spans the 20th to 80th percentiles of the
healthy bank. The random seed is fixed so that the sample is
reproducible.

The intent is to audit items that the pipeline itself considers
well-behaved and typical. There are two reasons for this. First, if the
gold standards fail on well-behaved items, then they fail everywhere.
Second, a sample drawn from the extremes would confound item pathology
with authoring error, since an item can be badly calibrated and correctly
written at the same time.

One property of the sample needs to be stated plainly, because it is a
limitation of the generator rather than of the sampling. The verdict
balance is 38 \textsc{violation} to 12 \textsc{compliant}. A balanced
split was not achievable. The Red Herring plus subtle violation mandate
of Section~\ref{sec:itemgen} makes hard items violation-bearing by
construction, and only 16 compliant-verdict items exist in the entire
quality-filtered hard-domain pool. This is a design lesson for any
regeneration of the bank, because a difficulty mandate that always plants
a violation teaches an examinee that the answer is always ``violation''.
I flag it here rather than presenting the sample as balanced.

\subsection{Two-step measurement}
Each reviewer works through the items in a fixed order and performs two
operations per item. The order of the two operations matters.

\begin{enumerate}[nosep,leftmargin=1.6em]
  \item \textbf{Independent verdict.} Reading only the scenario and the
        question, with the reference answer withheld, the reviewer
        records \textsc{violation}, \textsc{compliant}, or
        \textsc{ambiguous}.
  \item \textbf{Gold-standard grade.} The reference answer is then
        revealed and marked correct or incorrect. Correct requires both
        the right conclusion and a correct section-level citation. A
        written justification is required for every negative mark.
\end{enumerate}

Withholding the reference answer during the first step is what makes the
two measurements independent. If the order were reversed, the reviewer
would be anchored to the model's conclusion and the agreement rate would
be inflated. Reviewers are asked not to consult AI tools.
\textsc{Ambiguous} is offered as a first-class response, because an item
that two qualified professionals cannot resolve is a defective item, and
recording that is more useful than forcing a call.

\subsection{What the labels will estimate}
With two or more independent reviewers, the protocol yields three
quantities that the rest of this work cannot obtain.

\begin{itemize}[nosep,leftmargin=1.4em]
  \item \textbf{Gold-standard error rate}, which is the proportion of
        reference answers judged incorrect. This is a direct estimate of
        how much authoring error the closed loop has injected into the
        bank, and it is the single number most needed in order to
        interpret every $b$ reported here.
  \item \textbf{Inter-rater agreement}, measured by Cohen's $\kappa$ on
        the first-step verdicts. This calibrates the ceiling. If experts
        disagree with each other on these scenarios, then no automated
        judge can be expected to do better, and the difficulty scale is
        measuring genuine regulatory ambiguity as well as model ability.
  \item \textbf{Human and machine concordance}, which is the agreement
        between expert verdicts and the pipeline's verdicts, computed per
        domain. This indicates which domains need regeneration rather
        than recalibration.
\end{itemize}

\subsection{Pre-committing to the failure modes}
Two outcomes would materially damage the claims in this paper, and I
state them now rather than after seeing the labels. First, a
gold-standard error rate above roughly $10\%$ would mean the difficulty
estimates are contaminated at a level that recalibration cannot repair,
and the affected domains would need to be regenerated and recalibrated
rather than adjusted. Second, low inter-rater agreement ($\kappa < 0.4$)
would indicate that the items are ambiguous rather than hard, which would
mean the bank measures the resolution of under-specified scenarios rather
than regulatory knowledge. Either finding would be reported. Nothing in
the current results anticipates either one, and nothing in them excludes
either one.
